% Consolidated 32-moving-optics from 620defd2149c8dcfd52d29c147aa4708c27ab8f3; original TeX fragments preserved.

% === BEGIN 32-moving-optics-224-245.tex blob d5a82fc5844b01d51ce97d41d5b21eecad0bcbab ===
% Irodov Part V §5.6 Optics of Moving Sources, original problems 224--245.
% Source: Arihant PDF printed pp.224--226, answers pp.336--337.
% All frequencies in this section are measured in the observer frame unless
% marked proper; beta>0 denotes motion toward the observer when explicitly stated.

\movingopticsproblem{224}{224}
在斐索齿轮测光速实验中，齿轮距反射镜
$l=7.0\,\mathrm{km}$，齿数 $z=720$。
相邻两次光束消失对应转速
$n_1=283\,\mathrm{s^{-1}}$ 和
$n_2=313\,\mathrm{s^{-1}}$。
求真空光速。
\exsolution
光往返齿轮与反射镜所用时间为
$t=2l/c$。相邻两次消光意味着齿轮在这段时间里
多转过一个齿的\emph{角位置周期}，
即 $\Delta n\,t=1/z$，
其中 $\Delta n=n_2-n_1$。
于是
\[
 \boxed{c=2lz(n_2-n_1)
 =2(7.0\times10^3)(720)(30)
 \simeq3.02\times10^8\,\mathrm{m/s}}.
\]
\excheck 齿轮开口与齿距交替出现，
从一消光转速到\emph{下一次}消光转速的变化，
与同相齿位周期相联系；数据给出近似 $3.0\times10^8$。

\movingopticsproblem{225}{224}
一光源以相对于接收器的速度 $v\ll c$ 运动。
证明其发光频率的一阶相对多普勒变化
$\Delta\omega/\omega_0=(v/c)\cos\theta$，
其中 $\theta$ 是光源速度与\emph{指向观察者}的发射方向夹角。
\exsolution
设相邻两个波峰在光源系发射时间间隔为
$T_0=2\pi/\omega_0$。
在一阶近似下时间膨胀可忽略，
两次发射间光源朝接收方向移动
$vT_0\cos\theta$。
相邻波峰的距离因而成为
$\lambda\simeq cT_0-vT_0\cos\theta$，
所以
\[
 \omega_{\rm obs}=\frac{2\pi c}{\lambda}
 =\frac{\omega_0}{1-(v/c)\cos\theta}
 \simeq\omega_0\left[1+\frac vc\cos\theta\right].
\]
故
\[
 \boxed{\frac{\Delta\omega}{\omega_0}
 \simeq\frac vc\cos\theta}.
\]
\excheck 朝接收器运动对应 $\cos\theta>0$，发生蓝移；
一阶展开误差为 $O(v^2/c^2)$。

\movingopticsproblem{226}{224}
一束具有动能 $T=10\,\mathrm{MeV}$ 的激发态
$\mathrm{He}^+$ 离子发出本征波长
$\lambda_0=410\,\mathrm{nm}$ 的谱线。
沿与离子速度成 $\theta=30^\circ$ 的方向观察，
求多普勒波长移动。
\exsolution
离子静质量能
$Mc^2\simeq4u\,c^2\simeq3727\,\mathrm{MeV}$。
因 $T/(Mc^2)\ll1$，用
\[
 \beta=\frac vc\simeq\sqrt{\frac{2T}{Mc^2}}
 \simeq0.0733.
\]
向前侧的观察方向有蓝移，
线性近似
$\Delta\lambda\simeq-\lambda_0\beta\cos\theta$。
故
\[
 \boxed{\Delta\lambda\simeq
 -(410\,\mathrm{nm})(0.0733)\cos30^\circ
 \simeq-26\,\mathrm{nm}}.
\]
\excheck 由精确相对论多普勒公式
$\lambda_{\rm obs}=\lambda_0\gamma(1-\beta\cos\theta)$
可给出略不同的高阶修正；书末答案使用低速一阶值。

\movingopticsproblem{227}{224}
在太阳赤道两侧边缘观测
$\lambda_0=0.59\,\mu\mathrm m$ 谱线，
两边缘测得波长相差
$\delta\lambda=8.0\,\mathrm{pm}$。
求太阳的自转周期。
\exsolution
假设赤道接近侧面观察，
太阳两个对称边缘分别具有视向速度
$\pm v$，小速度近似波长差为
$\delta\lambda\simeq2(v/c)\lambda_0$。
故 $v=c\,\delta\lambda/(2\lambda_0)$。
若太阳半径为 $R_\odot$，
自转周期
\[
 P=\frac{2\pi R_\odot}{v}
 =\boxed{\frac{4\pi R_\odot\lambda_0}
 {c\,\delta\lambda}}.
\]
取 $R_\odot\simeq6.96\times10^8\,\mathrm m$，
\[
 \boxed{P\simeq2.15\times10^6\,\mathrm s
 \simeq25\,\mathrm d}.
\]
\excheck 这里用赤道边缘投影速度，
并忽略太阳自转轴的视线倾斜和差分自转。

\movingopticsproblem{228}{224}
一对质量相等、远到望远镜无法分辨的双星围绕共同
质心作圆周运动。观测到两颗星的谱线最大相对分裂
$(\Delta\lambda/\lambda)_{\max}
=1.2\times10^{-4}$，
相邻两次出现最大分裂的间隔为
$\tau=30\,\mathrm d$。
求双星间距离 $d$ 和每颗恒星的质量 $m$。
\exsolution
设轨道平面恰包含视线方向，
每星轨道速度为 $v$。
两星视向速度相反，最大相对分裂为
$\epsilon=2v/c$，
所以 $v=\epsilon c/2$。
一轨道周期内出现两次最大分裂，故轨道周期
$P=2\tau$，角速度
$\Omega=2\pi/P=\pi/\tau$。
每星绕质心的轨道半径为 $d/2$，
$v=\Omega d/2$，得到
\[
 \boxed{d=\frac{\epsilon c\tau}{\pi}
 \simeq2.97\times10^{10}\,\mathrm m
 \simeq3.0\times10^7\,\mathrm{km}}.
\]
由万有引力提供向心力
$Gm^2/d^2=m v^2/(d/2)$，
有 $m=2v^2d/G$，从而
\[
 \boxed{m=\frac{\epsilon^3c^3\tau}{2\pi G}
 \simeq2.9\times10^{29}\,\mathrm{kg}}.
\]
\excheck $P=2\tau$ 是本题关键；
若轨道倾角未知，上述速度、距离、质量是取
边缘正视轨道平面时的估计，质量将有 $\sin^{-3}i$
修正，须显式说明。

\movingopticsproblem{229}{224}
频率为 $\omega_0$ 的平面光波垂直入射一面
以相对论速度 $V$ 向入射光接近的反射镜。
求反射光在实验室系中的频率，
并给出 $V\ll c$ 的近似。
\exsolution
设 $\beta=V/c>0$。
在反射镜静止系，入射光经洛伦兹频率变换，
频率为
$\omega_{\rm in}'=\gamma(1+\beta)\omega_0$。
理想静止反射镜使光反向传播但不改变频率，
故 $\omega_{\rm out}'=\omega_{\rm in}'$。
变换回实验室系，再乘
$\gamma(1+\beta)$，
得到
\[
 \omega_{\rm refl}
 =\gamma^2(1+\beta)^2\omega_0
 =\boxed{\omega_0\frac{1+\beta}{1-\beta}}.
\]
当 $\beta\ll1$，
\[
 \boxed{\omega_{\rm refl}\simeq\omega_0(1+2V/c)}.
\]
\excheck 两次多普勒变换使一阶相对频移为
$2V/c$，而非 $V/c$。

\movingopticsproblem{230}{224}
雷达工作波长 $\lambda=50.0\,\mathrm{cm}$。
雷达发射波与迎面飞来飞机的反射回波之间，
产生拍频 $\Delta\nu=1.00\,\mathrm{kHz}$。
求飞机速度。
\exsolution
往返反射多普勒在 $v\ll c$ 时给
$\Delta\nu/\nu_0\simeq2v/c$，
而 $\nu_0=c/\lambda$，故
\[
 \boxed{v=\frac{\lambda\Delta\nu}{2}
 =\frac{(0.500)(1000)}2
 =250\,\mathrm{m/s}
 \simeq900\,\mathrm{km/h}}.
\]
\excheck 飞机接近时回波频率较高；
拍频只给速度大小，方向需另由频率升降判定。

\movingopticsproblem{231}{224}
利用平面波相位
$\varphi=\omega t-kx$ 的洛伦兹不变性，
推导沿 $+x$ 方向相对速度 $V$ 的惯性系
$K'$ 中频率与波数的变换式。限定一维传播。
\exsolution
坐标变换
$t'=\gamma(t-Vx/c^2)$，
$x'=\gamma(x-Vt)$，
$\gamma=(1-\beta^2)^{-1/2}$。
由同一相位在两系相等
$\omega t-kx=\omega't'-k'x'$，
收集 $t,x$ 的系数，得到
\[
 \omega=\gamma(\omega'+Vk'),\qquad
 k=\gamma\left(k'+\frac{V\omega'}{c^2}\right).
\]
解逆变换则为
\[
 \boxed{\omega'=\gamma(\omega-Vk)},\qquad
 \boxed{k'=\gamma\left(k-\frac{V\omega}{c^2}\right)}.
\]
对于沿 $+x$ 传播的真空光 $k=\omega/c$，
\[
 \boxed{\omega'=\gamma(1-\beta)\omega
 =\omega\sqrt{\frac{1-\beta}{1+\beta}}}.
\]
\excheck 波四矢量
$(\omega/c,k)$ 的 Minkowski 范数
$(\omega/c)^2-k^2$ 保持不变，
真空光的该范数在两系均为0。

\movingopticsproblem{232}{224}
星云中的氢谱线本征波长为
$\lambda_0=434\,\mathrm{nm}$，
观测到它向红端移动
$\Delta\lambda=130\,\mathrm{nm}$。
求星云的视向退行速度。
\exsolution
设观测波长为
$\lambda=\lambda_0+\Delta\lambda=564\,\mathrm{nm}$。
相对论视向多普勒公式为
\[
 \frac{\lambda}{\lambda_0}
 =\sqrt{\frac{1+\beta}{1-\beta}},
 \qquad \beta=v/c.
\]
平方解出
\[
 \boxed{\beta=
 \frac{(\lambda/\lambda_0)^2-1}
 {(\lambda/\lambda_0)^2+1}}
 =\frac{564^2-434^2}{564^2+434^2}
 \simeq0.256.
\]
故 $\boxed{v\simeq0.26c}$。
\excheck 红移对应星云远离，必须有 $\beta>0$；
速度已非极低速，不能仅使用 $v/c=\Delta\lambda/\lambda_0$。

\movingopticsproblem{233}{225}
汽车驾驶员应以多大速度朝向红色交通灯运动，
才能将发光波长约
$\lambda_0=0.70\,\mu\mathrm m$ 的红灯
看成波长约
$\lambda'=0.55\,\mu\mathrm m$ 的绿光？
\exsolution
观察者迎面运动，频率蓝移，
波长比为
$\lambda'/\lambda_0=\sqrt{(1-\beta)/(1+\beta)}$。
因此
\[
 \boxed{\beta=
 \frac{(\lambda_0/\lambda')^2-1}
 {(\lambda_0/\lambda')^2+1}
 =\frac{0.70^2-0.55^2}{0.70^2+0.55^2}
 \simeq0.236},
\]
从而
\[
 \boxed{v\simeq7.1\times10^7\,\mathrm{m/s}
 =7.1\times10^4\,\mathrm{km/s}}.
\]
\excheck 这是相对论理想观察色移，
而非现实道路上可能达到的驾驶速度。

\movingopticsproblem{234}{225}
观察者以 $v_1=c/2$ 沿直线运动，
其前方的单色光源也以
$v_2=3c/4$ 沿同一方向运动。
光源固有角频率为 $\omega_0$。
求后方观察者接收的角频率。
\exsolution
光源相对地面向前运动，朝后方发出光，
所以地面系接收的波频为
\[
 \omega_{\rm lab}=\omega_0
 \sqrt{\frac{1-\beta_2}{1+\beta_2}}
 =\frac{\omega_0}{\sqrt7}.
\]
后方观察者朝前运动，迎面遇到从前方传来的
后向传播光，因此在观察者系频率再次乘
$\sqrt{(1+\beta_1)/(1-\beta_1)}=\sqrt3$。
于是
\[
 \boxed{\omega=\omega_0\sqrt{\frac37}}.
\]
\excheck 两次变换应严格考虑光子的传播方向；
仅把两速度数值相减而不做相对论合成是不严谨的。

\movingopticsproblem{235}{225}
动能 $T=1.0\,\mathrm{MeV}$ 的一束激发态氢原子，
发射本征波长
$\lambda_0=656.3\,\mathrm{nm}$ 的谱线。
若沿垂直原子束的方向观察，
求横向多普勒位移。
\exsolution
横向多普勒效应来自时间膨胀，
在实验室观察方向与运动方向正交时，
观测波长为
$\lambda=\gamma\lambda_0$，
其中 $\gamma=1+T/(m_{\rm H}c^2)$。
因此
\[
 \boxed{\Delta\lambda=
 \lambda_0(\gamma-1)
 =\lambda_0\frac{T}{m_{\rm H}c^2}}.
\]
以 $m_{\rm H}c^2\simeq938\,\mathrm{MeV}$，
\[
 \boxed{\Delta\lambda\simeq0.700\,\mathrm{nm}}.
\]
\excheck 为红移；尽管横向视向速度为零，
二阶 $v^2/c^2$ 效应仍存在。

\movingopticsproblem{236}{225}
一个单色信号源沿直线以
$v=0.80c$ 匀速运动，固有角频率
$\omega_0=3.0\times10^{10}\,\mathrm{s^{-1}}$。
静止接收者位于离光源轨迹最近距离为 $l$ 的点 $P$
（原图5.37），$O$ 是最近接近点。
求接收频率：(a) 光源\emph{实际上刚到} $O$ 的时刻；
(b) 接收者\emph{刚看到光源处于} $O$ 的时刻。
\exsolution
设光源沿 $+x$，在 $t=0$ 经过 $O=(0,0)$，
接收者 $P=(0,l)$。
对被接收到的信号，发射时刻
$t_e<0$ 与光行程满足
$c(0-t_e)=\sqrt{l^2+(vt_e)^2}$，
由此
$t_e=-l/[c\sqrt{1-\beta^2}]$。
在发射时刻，从源指向接收者的光线
与运动方向的方向余弦恰为 $\cos\theta=\beta$。
相对论多普勒式
\[
 \omega=\frac{\omega_0}{\gamma(1-\beta\cos\theta)}
\]
于是 (a)
\[
 \boxed{\omega_a=\gamma\omega_0
 =\frac{\omega_0}{\sqrt{1-\beta^2}}
 =5.0\times10^{10}\,\mathrm{s^{-1}}}.
\]
(b) 当观察者收到的是\emph{从 $O$ 发出的光}时，
那束光的实验室发射角为 $\theta=90^\circ$，
故
\[
 \boxed{\omega_b=\frac{\omega_0}{\gamma}
 =\omega_0\sqrt{1-\beta^2}
 =1.8\times10^{10}\,\mathrm{s^{-1}}}.
\]
\excheck 必须区分光源现在的位置与光源的\emph{视位置}；
混淆发射时刻和接收时刻会将(a)(b)的数值弄反。

\movingopticsproblem{237}{225}
速度接近光速的一窄束电子沿周期
$d=2.0\,\mu\mathrm m$ 的金属光栅表面掠过，
电子运动方向垂直于栅线。
沿观察角 $\theta$ 看到彩色带状辐射（原图5.38）。
解释该现象，并求 $\theta=45^\circ$、
$v\simeq c$ 时的特征辐射波长。
\exsolution
电子与其在金属中的诱导镜像电荷形成
运动电偶极子；表面周期结构使电磁场按光栅周期受激。
相邻格点辐射在观察方向的相位差必须为整周，
得到 Smith--Purcell 辐射关系
\[
 \boxed{\lambda_m=\frac d m
 \left(\frac c v-\cos\theta\right)},\quad
 m=1,2,\ldots
\]
对题目要求的第一级光谱 $m=1$，
$v\simeq c,\theta=45^\circ$，故
\[
 \boxed{\lambda\simeq
 2.0\,\mu\mathrm m(1-\cos45^\circ)
 \simeq0.586\,\mu\mathrm m\approx0.6\,\mu\mathrm m}.
\]
\excheck $\lambda$ 随观察角变化，因此看见色带；
高次级次的波长是上式的 $1/m$，
本题书末取基本级次。

\movingopticsproblem{238}{225}
质量为 $m$ 的气体原子在温度 $T$ 下达到热平衡，
其固有发射角频率为 $\omega_0$。
证明：(a) 沿给定观察方向的多普勒展宽谱线满足
\[
 I_\omega=I_0\exp[-a(1-\omega/\omega_0)^2],
 \qquad a=\frac{mc^2}{2k_BT};
\]
(b) 求光强降至峰值一半时，
两个频率之间的相对全宽 $\Delta\omega/\omega_0$。
\exsolution
原子沿视线的速度分量
$v_\parallel$ 服从一维麦克斯韦分布：
\[
 f(v_\parallel)\,dv_\parallel=
 \sqrt{\frac{m}{2\pi k_BT}}
 e^{-mv_\parallel^2/(2k_BT)}\,dv_\parallel.
\]
对 $|v_\parallel|\ll c$，
$\omega\simeq\omega_0(1+v_\parallel/c)$，
即 $v_\parallel=c(\omega/\omega_0-1)$。
因 $|dv_\parallel/d\omega|=c/\omega_0$ 为常数，
谱密度与该高斯分布成正比，
故得到给定的 $I_\omega$。
在半高位置
$a(1-\omega/\omega_0)^2=\ln2$，
两根分居 $\omega_0$ 两侧，因此全宽
\[
 \boxed{\frac{\Delta\omega}{\omega_0}
 =2\sqrt{\frac{2k_BT\ln2}{mc^2}}}.
\]
\excheck 全宽是两半高频率\emph{之差}，不是
单侧半宽；原子速度为非相对论速率时高斯近似成立。

\movingopticsproblem{239}{226}
一平面电磁波在折射率为 $n$ 的透明介质中传播。
介质沿传播方向相对于实验室惯性系以
$V\ll c$ 匀速运动。
求电磁波在实验室系中的传播速度。
\exsolution
在介质静止系的相速度为 $u'=c/n$。
由洛伦兹速度合成（同方向）
\[
 u=\frac{u'+V}{1+u'V/c^2}
 =\boxed{\frac{c/n+V}{1+V/(nc)}}.
\]
令 $\epsilon=V/(nc)\ll1$，
一阶展开得菲涅耳拖曳公式
\[
 \boxed{u\simeq\frac cn+
 V\left(1-\frac1{n^2}\right)}.
\]
\excheck 当 $n=1$，结果恒为 $c$，
这就能检查一阶项的\emph{减号}；
原书纯文本答案此处易误识别为加号。

\movingopticsproblem{240}{226}
地球公转使黄道面内的恒星位置发生周年光行差。
恒星周年视位置的总摆动角
$\Delta\theta=41''$。
求地球的轨道速度。
\exsolution
地球以横向速度 $v$ 运动时，
对远方恒星的视位置一阶偏移角幅度约
$\alpha=v/c$（弧度）。
地球公转半周后横向速度反向，
两极端位置总角距约为
$\Delta\theta=2v/c$。
故
\[
 \boxed{v=\frac{c\Delta\theta}{2}}.
\]
将 $41''$ 转弧度
$\Delta\theta=41\pi/(180\times3600)
\simeq1.99\times10^{-4}$，
得
\[
 \boxed{v\simeq2.98\times10^4\,\mathrm{m/s}
 \simeq30\,\mathrm{km/s}}.
\]
\excheck 41角秒是\emph{峰峰值}，
若直接按单侧角移会多算一倍。

\movingopticsproblem{241}{226}
证明：对沿 $x$ 方向以速度 $V$ 作洛伦兹变换的
两惯性系 $K\to K'$，
光传播方向与 $x$ 轴的夹角满足
\[
 \cos\theta'=\frac{\cos\theta-\beta}
 {1-\beta\cos\theta},
 \qquad \beta=V/c.
\]
\exsolution
真空光速度在 $K$ 中的分量为
$u_x=c\cos\theta$，
$u_y=c\sin\theta$。
洛伦兹速度合成给
\[
 u_x'=\frac{u_x-V}{1-Vu_x/c^2}
 =c\frac{\cos\theta-\beta}
 {1-\beta\cos\theta}.
\]
由于光在 $K'$ 中仍满足
$|\bm u'|=c$，故 $u_x'=c\cos\theta'$，
于是
\[
 \boxed{\cos\theta'=
 \frac{\cos\theta-\beta}{1-\beta\cos\theta}}.
\]
\excheck 在 $\beta=0$ 时 $\theta'=\theta$；
在 $\theta=0$ 时光仍沿正 $x$ 传播。

\movingopticsproblem{242}{226}
观察者以 $V=0.99c$ 相对于地球运动。
利用上一题公式，求地球观察者所见的\emph{前半天球}
中全部恒星，在高速观察者视野中集中形成的圆锥
的半顶角。
\exsolution
注意看到的恒星方向与接收光子的运动方向相反。
设地球系中位于前半天球边缘的恒星，
其入射光子方向垂直于运动方向，即
$\cos\theta=0$。
由光行差公式，
$\cos\theta'=-\beta$。
因此恒星\emph{视线}方向与 $+x$ 运动方向的夹角
$\alpha$ 满足
$\cos\alpha=-\cos\theta'=\beta$，
于是
\[
 \boxed{\alpha=\arccos0.99\simeq8.1^\circ}.
\]
\excheck 这里给的是圆锥半顶角（视线与轴之夹角）；
若将“全顶角”定义为两侧母线夹角，则为约 $16.2^\circ$。

\movingopticsproblem{243}{226}
带电粒子以恒速 $v$ 穿过折射率为 $n$ 的均匀透明媒质。
从不同轨迹点发出的电磁波干涉，推导
发生瓦维洛夫--切伦科夫辐射的速度条件与圆锥方向。
\exsolution
设粒子在相隔时间 $\Delta t$ 的两个位置
$A,B$ 相距 $v\Delta t$。
粒子到达 $B$ 时，$A$ 点此前激发的光波
已经传播距离 $c\Delta t/n$。
要使两处发射波在与轨道成 $\theta$ 的远场方向
同相叠加，需沿辐射方向的路径差补偿传播时差：
\[
 v\Delta t\cos\theta=\frac cn\Delta t.
\]
故
\[
 \boxed{\cos\theta=\frac c{nv}},
 \qquad \boxed{v>\frac cn}.
\]
只有 $nv/c>1$ 才允许非零开角的相干辐射，
其光线沿以粒子速度方向为轴、
半顶角 $\theta=\arccos[c/(nv)]$ 的圆锥传播。
\excheck $v<c$ 仍可能满足 $v>c/n$，
因此无须粒子超光速；这里的阈值是介质相速度。

\movingopticsproblem{244}{226}
折射率为 $n=1.60$ 的介质中，
求电子与质子产生切伦科夫辐射所需的最小动能。
再问：哪种粒子在此介质中的阈值动能为
$T_{\min}=29.6\,\mathrm{MeV}$？
\exsolution
阈值速度为
$\beta_{\rm th}=1/n$。
对应相对论因子
\[
 \gamma_{\rm th}
 =\frac1{\sqrt{1-1/n^2}}
 =\frac n{\sqrt{n^2-1}},
\]
故任一静质量 $m$ 的带电粒子阈值动能为
\[
 \boxed{T_{\min}=
 \left(\frac n{\sqrt{n^2-1}}-1\right)mc^2}.
\]
取 $n=1.60$，$\gamma_{\rm th}-1\simeq0.281$。
电子 $m_ec^2=0.511\,\mathrm{MeV}$，
质子 $m_pc^2=938\,\mathrm{MeV}$，得
\[
 \boxed{T_{e,\min}\simeq0.144\,\mathrm{MeV}},\qquad
 \boxed{T_{p,\min}\simeq264\,\mathrm{MeV}}.
\]
由 $29.6/0.281\simeq105\,\mathrm{MeV}$
推回静质量能，符合 $\mu$ 子（缪子）。
\excheck 同一介质对不同粒子的阈值
$\beta_{\rm th}$ 相同，但动能正比于静质量能。

\movingopticsproblem{245}{226}
电子在折射率 $n=1.50$ 的介质内
发出与运动方向成 $\theta=30^\circ$
的切伦科夫光。
求电子的动能。
\exsolution
切伦科夫角条件
$\cos\theta=1/(n\beta)$，故
\[
 \boxed{\beta=\frac1{n\cos\theta}},
 \quad
 \boxed{\gamma=
 \frac{n\cos\theta}{\sqrt{n^2\cos^2\theta-1}}}.
\]
代入
$n=1.50,\theta=30^\circ$，
$n\cos\theta=1.2990$，
故 $\beta\simeq0.7698$、$\gamma\simeq1.567$。
于是
\[
 \boxed{T=(\gamma-1)m_ec^2
 =\left(\frac{n\cos\theta}
 {\sqrt{n^2\cos^2\theta-1}}-1\right)
 m_ec^2
 \simeq0.290\,\mathrm{MeV}}.
\]
\excheck 按给定精确公式和电子静能 $0.511\,\mathrm{MeV}$
直接计算得到约 $0.29$MeV；
本版书末若印 $0.23$MeV，与这些参数不相容，
已将这项数值差异标入勘误以待逐版校对。
% === END 32-moving-optics-224-245.tex ===
